\scanpage{81}
% archive sheet number: 80
% The typed draft is cancelled by oblique strokes across the page. Marginal pen additions are retained separately.
\deleted{et l'hypothèse de récurrence montre alors qu'on a l'inégalité}
\[
\deleted{\operatorname{long}M/tM\qquad\operatorname{long}M_{\underline{p}_i}\ ,}
\]
\deleted{l'égalité étant vraie sss les égalités analogues sont vraies pour $M'$ et $M''$, ce qui pqr l'hypothèse de récurrence signifie précisément que b) et c) sont vérifiés\inserted{pour $M$}. CQFD}

\deleted{\uline{Corollaire} 3. Sous les conditions de 1.\inserted{(ii)}, supposons que $\operatorname{long}(F/tF)_z$ soit égal à 1, alors il existe une seule composante irréductible $S_i$ de $S=\operatorname{supp}F$ passant par $Z$, et si $x$ est son point générique, on a $\operatorname{long}F_x=1$. De plus, $z$ est un point régulier de $\deleted{S_i}\inserted{\unsure{S_{\mathrm{réd}'}|S_{\mathrm{réd}},}}$ et $t$ engendre l'idéal maximal de $\underline{O}_{S_{\mathrm{réd}},z}$.}

\deleted{\uline{Proposition} 4\quad Sous les conditions de 1., soit $x\quad X$, supposons que $V(t)$ ne contienne aucune composante irréductible $S_i$ de $S=\operatorname{supp}F$ passant par $x$, et que $F/tF$ soit réduit en $x$. Alors $t$ est $F$-régulier en $x$, et $F$ est réduit en $x$.}

\deleted{Pour la première assertion, la démonstration est donnée dans 5.6.2. Le fait que $F_x$ n'ait pas de cycles premiers associés immergés passant par $x$ est \deleted{un cas} alors un cas particulier de \unsure{8|6}. ci-dessous. Reste à prouver que les long\inserted{u}eurs de $F$ aux points générisations maximales de $x$ dans $S$ sont égales à 1, ce qui est contenu dans le corollaire 3 précédent; on pourrait d'ailleurs ajouter que les générisations maximales de $x$ dans $\operatorname{supp}F/tF$ sont réguliers sur $S_{\mathrm{réd}}$ et que $t$ \unsure{engendre} l'idéal maximal des anneaux locaux de $S_{\mathrm{réd}}$ en ces points \ldots}

\deleted{\uline{Proposition} 5. Sous les conditions de 1., soit $T$ un cycle premier associé à $F$ (pas nécessairement maximal), $T'$ une composante de $T\quad V(t)$, \deleted{\unsure{supposons} de point générique $x$.} Supposons que $t$ est $F$-régulier au point générique $x$ de $T'$, alors \deleted{$T'$ est un} $x\quad\operatorname{Ass}F/tF$.}
\margin{\handwritten{$\left.\begin{gathered}\text{\unsure{Corollaire|Condition} pour}\\\unsure{x|t}\ \text{\unsure{soit|est}}\ F\ \text{\unsure{régulier}.}\end{gathered}\right]$}}

\deleted{La démonstration figure dans 8.1.4.}

\deleted{\uline{Corollaire} 6. Sous les conditions de 1., soit $x\quad X$, \deleted{et} supposons que $t$ soit $F$-régulier en $x$ et que $F/tF$ n'ait pas de cycles premiers associés immergés passant par $x$. Alors \inserted{\handwritten{$F$ n'a pas de cycles premiers immergés \unsure{passant} par $x$.}} les composantes irréductibles de $\operatorname{supp}F/tF$ passant par $x$ sont \unsure{e xactement|exactement} les composantes irréductibles passant par $X$ des $S_i\cap V(t)$, où les $S_i$ sont les comp irred de $\operatorname{supp}F=S$}
